\chapter{离子键：电子转移之后发生什么}

\begin{kaichang}
去厨房找出食盐，倒一小撮在白纸上，用放大镜看一看——或者用手机相机放大拍一张照片。你会发现：盐粒不是圆圆的沙粒，而是一颗颗方方正正的小晶体，像被缩小的骰子，棱角分明。找一颗大些的盐粒，垫着纸用锤子轻轻敲碎，再看碎片：还是方的，只是更小。

现在请你猜三件事。第一猜：盐为什么这么“方”？方得这么整齐，是谁规定的？第二猜：盐一遇水就消失得无影无踪，可把它撒在炉灶的火焰上，它却只是让火焰变黄（这是钠的焰色，第 9 章讲过），自己迟迟不肯熔化——既怕水又不怕火，这是为什么？第三猜：干燥的盐完全不导电，把同一勺盐溶进一杯水里，这杯水却能导电了。盐没变多也没变少，发生了什么？

把三个猜测写在纸上。这一章结束时，我们要用盐粒内部那座整齐得惊人的“城市”来一一对答案。
\end{kaichang}

\section{先把盐的脾气摸清楚}

动手之前，先把看得见的事实摆齐。取食盐、白糖各一小勺，逐项对比，你会发现一堆耐人寻味的差异：

\begin{itemize}[itemsep=2pt]
\item \textbf{外形}：盐是立方体小晶体，糖的颗粒也有棱角，但敲碎后形状不规则；
\item \textbf{硬度}：两者用手指都捏不碎，都挺硬；
\item \textbf{抗打击}：锤子一敲，糖碎成乱碴，盐却裂成边缘整齐的小方块，像沿着隐形的缝线拆开；
\item \textbf{耐火}：糖加热到一百多摄氏度就熔化、变黄、变焦；盐在普通火焰上烤很久也只是发红，纹丝不动——纯盐要烧到 $801\,^{\circ}\mathrm{C}$ 才熔化；
\item \textbf{亲水}：两者都易溶于水，但干燥时是绝缘体，水溶液的行为却大不相同（这一点先记下，后面做实验见分晓）。
\end{itemize}

同一种白色小颗粒，脾气处处不同。糖熔化变焦时的那些事属于另一类键的故事（第 19 章），本章只盯住盐：它的方、它的脆、它的耐火、它“下水前后”导电性的变脸，全部指向同一个解释——盐不是一种“分子”，而是一种\textbf{带电粒子的巨大排列}。

\section{电子移交之后：两个离子的诞生}

先把镜头推进到一颗盐粒内部，看看组成它的两种粒子是怎么来的。这一幕在前面几章其实已经排练过：

钠原子有 11 个电子，按楼层排成 2、8、1——最外层孤零零 1 个电子，离核远、被内层电子遮挡得严，很容易被拽走（第 12 章讲过，这正是金属钠遇水暴躁的根源）。氯原子有 17 个电子，排成 2、8、7——最外层差 1 个住满，原子核对电子的“手劲”又大，是周期表上有名的抢电子高手（第 13 章）。用第 11 章的比分牌说：钠的电负性只有约 0.9，氯高达约 3.2，差距悬殊。

这样的两个原子相遇，结果没有悬念：钠最外层那个电子被氯整个夺过去。钠变成带一个正电荷的钠离子，电子层收缩成 2、8；氯变成带一个负电荷的氯离子，外层补成 2、8、8。

听起来公平交易、皆大欢喜？先别急着鼓掌。第 17 章立下的规矩是：一切变化都要算能量账。这笔账拆开算，前半截居然是\textbf{亏本}的——把电子从钠原子上硬拽下来，要倒贴能量（电离能，第 11 章）；电子投进氯原子怀里，能收回一部分（电子亲和能），但收不回全款。如果故事到此为止，钠和氯根本懒得反应。

真正让这笔交易大赚特赚的，是下一步：不计其数的正负离子，靠静电吸引紧紧地、整齐地挤到一起。理解这一步，就是理解离子键的全部。

\section{不是“一个对一个”，而是一座城市}

这里要纠正一个教科书无意中种下的印象。很多书写“一个钠离子和一个氯离子结合成一个氯化钠”，配上两个小圈挨在一起的图。这个画面错得相当离谱。

正负离子之间的吸引是\textbf{没有方向、没有名额限制}的：一个正离子不只吸引“它的”那个负离子，它吸引四面八方所有够得着的负离子；负离子也一样。于是最省能量的安排不是两两配对，而是一种三维的交替阵列——每个钠离子前后、左右、上下各被一个氯离子围住，共 6 个近邻；每个氯离子同样被 6 个钠离子围住。这样的阵列向四面八方延伸，直到铺满整颗盐粒：

\begin{center}
\begin{tikzpicture}[scale=0.55]
\foreach \x in {0,1,2,3,4}
  \foreach \y in {0,1,2,3,4} {
    \pgfmathparse{mod(\x+\y,2)}
    \pgfmathtruncatemacro{\parity}{\pgfmathresult}
    \ifnum\parity=0
      \fill[blue!60] (\x,\y) circle (0.18);
    \else
      \draw[red!70,fill=red!15] (\x,\y) circle (0.32);
    \fi
  }
\draw[gray!60] (0,0) grid (4,4);
\node[blue!80] at (2,-0.8) {\scriptsize 蓝点：钠离子\quad 红圈：氯离子（截面示意）};
\end{tikzpicture}
\end{center}

这是晶体一个截面的示意画法：圆的大小和颜色是人为的表示法，离子不是实心小球，真实的氯离子其实比钠离子大将近一倍（第 11 章量过它们的“腰围”）。关键信息只有一个：\textbf{异号相邻、同号相望，交替铺满}。

所以，一颗盐粒里不存在任何“氯化钠分子”。不存在哪一对离子彼此更亲、自成一户——每个离子都属于整座城市。夸张一点说：\textbf{一整颗盐粒就是一个“分子”}，一座由静电吸引粘合的离子之城。

这座城有多大？算一笔账。一颗边长 1 毫米的小盐粒，体积约 1 立方毫米，食盐的密度约 \ud{2.16}{g/cm^{3}}，所以质量约 $2.2\times 10^{-3}$ 克。氯化钠的摩尔质量约 \ud{58.4}{g/mol}，折合 $3.7\times 10^{-5}$ 摩尔，乘上阿伏伽德罗常数 \ud{6.02\times 10^{23}}{mol^{-1}}，约含 $2.2\times 10^{19}$ 个“钠氯单元”——也就是超过 $4\times 10^{19}$ 个离子。一小撮盐里的离子，比全地球人口每人再分十亿倍还多。

\section{规则层：静电的账本}

把这座城市粘住的力，你其实不陌生——就是摩擦气球能吸起头发的那一种：静电力。它的规律叫库仑定律：两个带电粒子之间的力，与两个电荷量的乘积成正比，与距离平方成反比。异号相吸，同号相斥。

这一条规律，加上“离子要尽量避免同号做邻居”，就同时解释了盐的所有怪脾气：

\textbf{为什么硬、熔点高？} 要让晶体熔化，就得让离子们挣脱四面八方所有邻居的吸引、开始自由滑动。每对近邻的吸引都不算惊人，但每个离子有 6 个近邻，近邻之外还有次近邻、再次近邻……整座城市加起来，是一笔巨大的“凝聚力存款”。要把 1 摩尔食盐晶体拆成彼此远离的气态离子，需要约 \ud{787}{kJ/mol}——这个量叫\textbf{晶格能}，是离子晶体牢不牢的度量。晶格能大，熔点就高：食盐 $801\,^{\circ}\mathrm{C}$ 才熔化，而糖分子之间没有这种静电粘合，一百多摄氏度就散架了。

\begin{qiaoliang}
现在可以兑现钠与氯相遇的能量总账了（每摩尔计）：
\begin{itemize}[itemsep=1pt]
\item 钠原子交出电子：$+496$（电离能，倒贴）；
\item 氯原子收下电子：$-349$（电子亲和能，回款）；
\item 气相中这一步小计：$+147$——\textbf{亏本}；
\item 离子聚成晶体：$-787$（晶格能，大赚）。
\end{itemize}
总账：$+147-787=-640$（kJ/mol），大降价，交易成交。请注意两件事。第一，反应的驱动力不是“氯想凑满 8 个电子”——原子没有愿望；是\textbf{整座城市}的能量比分开的原子低得多。第二，若只看前两项，钠和氯在气相里一对一反应是亏本的；是晶格能这笔“城建红利”扭转了全局。这就是为什么离子键从来以晶体形式出现，而不是以孤立离子对的形式出现。
\end{qiaoliang}

\textbf{为什么脆？} 硬和脆常常一起出现，原因却值得单独想。盐晶体之所以硬，是因为静电吸引强；之所以脆，是因为这种吸引对排列方式极其挑剔。想象沿某个晶面把上半块晶体推滑半个格距：原来异号相对的离子，现在变成同号脸对脸——吸引瞬间翻转为排斥，“咔”，晶体沿那个面干脆利落地裂开。这就是你敲盐粒时看到的整齐断面。金属被锤打时只会变形不会崩裂，靠的是完全不同的粘合方式（第 21 章细讲）。

\textbf{为什么有的离子晶体更“硬核”？} 库仑定律里有两个旋钮：电荷量和距离。电荷越大、离子越小（挤得越近），吸引越强，晶格能越大。比较一下：食盐里是 $+1$ 和 $-1$ 的离子，熔点 $801\,^{\circ}\mathrm{C}$；氧化镁里是 $+2$ 和 $-2$ 的离子，还都更小，晶格能约为食盐的 5 倍，熔点高达 $2852\,^{\circ}\mathrm{C}$——高到被用作炼钢炉的耐火砖。电荷密度（单位体积里塞进的电荷）是预判离子晶体脾气的好尺子。

\section{符号层：NaCl 是配方，不是分子}

现在才让符号正式登场，因为你已经知道它们该说什么、不该说什么。

钠与氯的反应写作：
\[\chem{2Na} + \chem{Cl_2} \rightarrow \chem{2NaCl}\]
两个生成物里的 \chem{NaCl}，读作“氯化钠”，但它只是\textbf{最简式}：它只承诺“在这座晶体里，钠离子和氯离子的数目之比是 $1:1$”，不承诺存在任何“\chem{NaCl} 分子”。这就像“这座城市男女比例 $1:1$”——是一句人口统计，不是说全城的人都是成对锁死的。

同样的逻辑读出别的离子化合物的配方：\chem{MgCl_2} 里镁离子带 $+2$、氯离子带 $-1$，所以一个镁离子配两个氯离子，正负电荷恰好抵消；\chem{Na_2O} 里两个钠离子配一个氧离子。离子晶体的化学式，本质是\textbf{电荷平衡表}：先数清每个离子带多少电，再凑出正负相抵的最小比例。碳酸钙 \chem{CaCO_3} 里那个 \chem{CO_3^{2-}} 是几个原子抱成一团的“原子团离子”——它内部原子之间用的是另一种键（第 19 章的共价键），但整体对外表现为一个带 $-2$ 电荷的离子，照样进城排队。

只有在食盐蒸气那种罕见场合，才存在短暂的“一个钠离子加一个氯离子”的孤立离子对——那才是勉强可以叫“\chem{NaCl} 分子”的东西，而你这辈子基本见不到它。

\section{证据层：X 射线“看见”了格子}

\begin{zhengju}
“食盐里没有分子”这个结论，不是靠聪明猜出来的，是被一束 X 射线硬逼着承认的。

故事从 1912 年开始。德国物理学家劳厄在想一个波的问题：当时 X 射线刚被发现十几年，没人确定它是不是一种波。劳厄推理：如果 X 射线是波长极短的波，而晶体内部果真是原子规则排列的点阵，那么 X 射线穿过晶体时，应该像光穿过密纹一样发生衍射，在底片上打出有规律的斑点。实验一做，底片上果然出现整齐的点阵图样——一石二鸟：X 射线是波，晶体是点阵。

第二年，英国的布拉格父子——父亲威廉·亨利·布拉格和儿子劳伦斯·布拉格——把这个现象变成了“量晶体内部间距的尺子”。衍射斑点的位置和角度，可以反推出点阵里相邻粒子的距离。他们选的头一批测量对象里，就有食盐。

测量结果让化学界很是别扭。大家本来期待在盐晶体里找到一对一对挨得特别近的“氯化钠分子”——一个钠配一个氯，像舞会上成对的舞伴。结果数据说得很清楚：晶体里只有一种排列，钠和氯交替占位、间距处处相同，没有任何两个粒子“格外亲近”。\textbf{期待中的分子对根本不存在。} 老布拉格起初也不甘心，试过别的排列假设，可只有“交替离子阵列”算出来的斑点位置能和底片严丝合缝。证据压倒习惯：化学家只好接受，食盐里没有食盐分子。

顺带地，这套方法还量出了离子的“腰围”，证实了另一件反直觉的事：钠离子比钠原子小将近一半，氯离子比氯原子大了一圈——和第 11 章从周期律推出的结论吻合。1915 年，布拉格父子同获诺贝尔物理学奖，小布拉格当年 25 岁，至今仍是诺贝尔科学奖最年轻的得主之一。这个故事值得记住的方法论是：他们手里没有任何一张照片能直接“看见”离子，只有一堆斑点；是把“如果排列是这样，斑点就该在那儿”的反推一遍遍做扎实，才让看不见的格子变成了公认的证据。
\end{zhengju}

\section{导电的变脸：不是有电，而是能搬家}

回到开场第三猜。干燥的盐明明由带电粒子组成，为什么不导电？

答案藏在一个容易忽略的区别里：导电需要的不是“存在电荷”，而是“\textbf{电荷能移动}”。电流就是电荷的定向流动。在固态盐里，每个离子都被邻居的吸引力锁死在格点上，只能原地哆嗦（热振动），哪儿也去不了——满腹电荷，动弹不得，自然是绝缘体。

两种办法给离子“解禁”：一是加热到熔化，晶格散开成液体，离子可以四处游走；二是溶进水里，水分子把离子从城墙上一个个请下来（怎么请，第 23 章细讲），离子在水里自由游泳。这两种状态下接上电源，正离子奔向负极、负离子奔向正极，两股相向的离子流就形成了电流。注意这和金属导电不是一回事：金属里跑动的是电子（第 21 章），盐溶液里跑动的是整个离子。

\begin{shiyan}
\textbf{三杯水的导电测试}

材料：9V 电池一块（或两节 5 号电池加电池盒）、小蜂鸣器一只（或 LED 串一只约 220 欧电阻）、带鳄鱼夹的导线三四根、食盐、白糖、纯净水、三个干净杯子、两根石墨电极（可用削尖的铅笔芯代替）。

步骤：① 把电池、蜂鸣器和两根电极串成一个回路，先把两根电极直接碰一下，确认蜂鸣器会响（回路是好的）。② 把两根电极插进一堆干食盐里，彼此不接触——蜂鸣器响不响？③ 把盐溶进半杯纯净水，电极再插进去（仍不互碰）——这次呢？④ 换一杯同样浓度的白糖水重复。⑤ 还可以凑近观察电极表面：盐水中电极上常会冒出细小气泡，那是电流在分解水（第 12 章的钠就是这么从熔盐里“电解”出来的，只是工业上用高温熔盐）。

观察点：干盐不响、盐水响、糖水基本不响——盐在固体时离子被锁住，入水后离子自由了；糖溶解后仍以分子形式存在，不产生可以搬家的电荷。

安全提示：只用电池，\textbf{绝对不许用插座里的市电}；两根电极不要直接相碰（会短路烧电池）；铅笔芯插水后会变脆，实验后丢弃；所有器具实验后洗净，不许尝任何实验用品。
\end{shiyan}

这个“解禁才导电”的性质不只是厨房趣闻，它是一座大工业的基石。把食盐加热到熔融再通强电流，就能迫使离子向两极搬家并在极板上交出、夺回电子——阴极析出银亮的金属钠，阳极冒出黄绿色氯气。第 12 章那颗遇水就炸的金属钠，就是这样从安安静静的白盐里“电解”出来的；电解食盐水溶液（氯碱工业）则同时产出烧碱、氯气和氢气，是化肥、塑料、消毒剂和造纸工业的原料起点。而医院里的生理盐水之所以能配进输液、海水之所以能导电、你神经细胞的电信号之所以能传导，靠的都是同一批自由移动的离子——后面讲生命时它们还会反复登场。

\section{有的盐溶于水，有的宁死不屈}

还剩最后一个问题：同样是离子晶体，食盐见水就化，可鸡蛋壳和石灰石（主要成分碳酸钙）泡在水里几年也不见少；医院做肠胃造影让病人喝的“钡餐”（硫酸钡），更是难溶到极点。区别在哪？

溶解是一场拔河。一头是晶格能：把离子从城墙上拆下来，要付“拆迁费”，城市粘得越牢，费用越高。另一头是水：水分子会把拆下来的离子团团围住、妥善安置，放出“安家费”（水合放热）。安家费盖过拆迁费，溶解就顺畅；拆迁费太高，晶体就赖着不走。食盐的离子电荷低、晶格能不算夸张，水付得起安家费，所以易溶（$25\,^{\circ}\mathrm{C}$ 时 100 克水能溶约 36 克）；碳酸钙里 $+2$ 对 $-2$，晶格能陡增，水就力不从心，每升水只能溶走十几毫克。拔河的细节——水分子凭什么围得住离子、温度和压强怎么搅局——第 23 章和第 25 章会专门清算，这里只要立住框架：\textbf{溶解度是“拆格子”与“安顿离子”两笔能量的差价}。

这个差价在医学上救人也伤人。钡离子 \chem{Ba^{2+}} 其实有毒，可溶的氯化钡足以致命；但硫酸钡难溶到喝下去几乎不放出钡离子，穿肠而过不吸收，又不透 X 射线——正好灌进肠胃当造影剂，让医生在 X 光片上看清消化道的轮廓。同一种元素，是毒药还是诊断工具，全看它锁在哪种晶体里。这是“剂量、形态与证据”原则的又一课：离开化合物形态谈“某元素有毒没毒”，是没有意义的。

冬天路面积雪时撒的融雪盐（常用氯化钙）也在用这套原理：氯化钙溶解时放出的热帮雪融化，溶出的离子还会干扰水结冰（第 25 章讲溶液性质时再算这笔账）。

\section{边界层：纯粹的离子键并不存在}

本章的模型——电子整个过户、离子像带电磁球一样整齐排列——对食盐、氧化镁这类“电负性差悬殊”的组合非常好用。但它有一条清晰的边界：当正离子电荷高、个头小，而负离子个头大、电子云蓬松时，正离子会把负离子的电子云明显拉向自己，电子不再是“彻底过户”，而变成了“重度偏向的共享”。这时晶体就开始不像典型的离子晶体：比如氯化铝，熔点低、受热容易直接升华，行为更像分子。

换句话说，\textbf{离子键和共价键不是两间互不相干的屋子，而是一个连续光谱的两端}。第 11 章的电负性差，就是粗略判断落在光谱哪一段的尺子：差越大越靠离子端，差越小越靠共享端。食盐恰好坐在离子的极端，所以本章的故事干净利落；但世界上绝大多数的键，都住在光谱的中间地带。

\begin{wuqu}
误解：“食盐溶于水，是一个个氯化钠分子被水冲散了。”

反例有两个。其一，本章的 X 射线证据：固体食盐里本来就没有“氯化钠分子”，只有离子的交替阵列，谈不上“分子被冲散”。其二，把食盐加热到熔化——全程不碰一滴水——熔融的盐照样导电，说明离子一旦脱离晶格就能自由移动、独立存在。溶解不是“拆分子”，而是“拆格子”：水分子把离子从城墙上请下来，分别围护着带走。这个区别不是抠字眼——它决定了为什么盐水能导电而糖水不能，也决定了你身体里的“盐”从来不是以“分子”的身份在工作。
\end{wuqu}

\section{回到那颗盐粒}

现在，兑现开场承诺。

盐为什么是方的？因为内部的离子点阵本身就是立方网格，晶体的外形是内部排列的宏观写真——亿万离子按立方队列站好，堆出来的外形自然是方方正正。敲碎后还是方的，因为敲击总是沿同号离子对脸的薄弱晶面劈开，碎块仍是小一号的立方队列。

为什么怕水不怕火？怕水，是因为水付得起安家费，把离子一个个请下城墙；不怕火，是因为把整座城市拆散要支付巨大的晶格能，普通火焰的能量差得远，要 $801\,^{\circ}\mathrm{C}$ 才能熔化。

为什么干盐不导电、盐水导电？因为导电的关键从来不是“有没有电荷”，而是“电荷能不能搬家”。固态时离子被静电锁在格点上，入水之后才恢复自由身。

你看，一颗最普通的盐粒，把粒子、规则、符号、证据串成了一整条链——这正是这本书一直在走的路。

\section{伏笔：电子一定要“整个过户”吗？}

钠和氯的故事之所以痛快，是因为双方实力悬殊：一个守不住，一个抢得狠，电子干脆利落整个过户。可世界上大多数原子相遇时没有这么大的差距。氢遇到氯，碳遇到氧，氧遇到氢——谁也没本事把对方的电子彻底夺走，僵持的结果只能是\textbf{共享}。

共享的电子怎样把原子拴在一起？为什么共享出来的东西才是“真正的分子”——有固定形状、有独立身份、能从水变成蒸气又变回来？共享如果不平均，又会引出什么后果？这些问题的答案，铺成了通往有机分子、水和生命的大门。下一章，第 19 章，共价键。

\begin{tiaozhan}
为什么晶格能比“一对离子的吸引”大那么多？可以把账逐项列出来：晶体里每个离子受到 6 个最近邻异号离子的吸引，又受到 12 个次近邻同号离子的排斥，再往外是 8 个异号、6 个同号……正负交替、随距离平方衰减。把这无穷多项加在一起，数学上收敛成一个纯数字——对食盐型结构约等于 $1.748$，叫\textbf{马德隆常数}。它的意思是：在晶体里，每个离子享受到的净静电吸引，是一对孤立离子对的约 $1.75$ 倍。城市确实比两人世界更省钱，而且可以省得精确算出来。把马德隆常数乘进库仑定律，再补上电子云太近时的排斥修正，就能从头算出食盐的晶格能，与实验值 \ud{787}{kJ/mol} 吻合到百分之几以内——这是“离子键就是静电力”这个模型最硬的定量胜利。用同样思路把本章能量账的每一步（升华、电离、亲和、成键）画成一张能量循环图，就是物理化学里有名的“玻恩–哈伯循环”。这一段跳过不影响主线，但若你想知道“库仑定律凭什么解释一整座晶体”，这里是入口。
\end{tiaozhan}

\section*{练一练}

\begin{sikaoti}
\item 画一幅 $4\times4$ 的二维离子阵列（仿照本章示意图），任选一个正离子：圈出它最近的几个异号离子和几个同号离子，用箭头标出吸引与排斥。用这幅图解释：晶体为什么既不会散架，也不会坍缩成一团？
\item 画一张能量台阶图：横轴为反应进程，纵轴为能量，依次画出“钠电离 $+496$”“氯得电子 $-349$”“离子聚成晶体 $-787$”（单位 kJ/mol）三步，标出净变化。再回答：如果有一种金属的电离能高得离谱，它还容易形成离子晶体吗？为什么？
\item 不查资料，从“电荷量”和“离子半径”两个旋钮出发，预测氧化镁 \chem{MgO} 的熔点比食盐高还是低，并说明理由；再预测氧化镁与氧化钙 \chem{CaO}（钙离子比镁离子大）谁的熔点更高。然后查数据核对你的推理。
\item 桌上有两杯无标签的透明液体，一杯是浓食盐水，一杯是浓糖水。用本章实验的器材设计一个鉴别方案，写出依据的原理。进一步想一想：如果某杯液体也能导电，能不能就此断定它的溶质在固态时一定是离子晶体？（提示：醋酸是分子，它的水溶液也能导电——它在水中“就地”产生了离子。第 26 章会讲这件事。）
\item 碳酸钙难溶于水，可鸡蛋壳泡进醋里会慢慢溶解并冒泡。用“拆格子的拆迁费 vs. 水的安家费”框架，提出一个解释：醋在这场拔河里可能帮了哪一边、怎么帮的？（提示：泡泡是二氧化碳，它从溶液里跑掉了。）
\item 本章算过，一颗边长 1 毫米的小盐粒约含 $4\times10^{19}$ 个离子。若把这些离子一个挨一个排成一列，每个离子占约 $0.3$ 纳米，这列队伍有多长？和地球到月球的距离（约 38 万千米）比一比。
\end{sikaoti}
